\section{The original arithmetic closed point lifts over the actual length-two Witt ring}

\begin{theorem}[An actual Witt point in the original principal open]
\label{thm:arithmetic-original-closed-point-witt-lift}
\lean{EnvelopeProperty.ArithmeticClosedPointLift.exists_large_characteristic_witt_lift}
\leanok
\uses{thm:arithmetic-large-characteristic-closed-point,def:witt-two-reduction,lem:witt-two-prime-annihilator,lem:witt-two-prime-kernel}
Let $R$ be a characteristic-zero integral domain of finite type over
$\mathbb Z$, smooth over $\mathbb Z$. Given the original nonzero
$r\in R$ and $N\ge0$, there are an original maximal ideal $\mathfrak m$,
a prime $p>N$, and a finite original residue field $\kappa=R/\mathfrak m$
of characteristic $p$, with $r\notin\mathfrak m$, and an actual ring map
\[
 \lambda:R\longrightarrow W_2(\kappa)
   =\operatorname{TruncatedWittVector}_{p,2}(\kappa)
\]
such that
\[
 \rho\circ\lambda=\varphi,
 \qquad \lambda(r)\in W_2(\kappa)^\times.
\]
Here $\varphi:R\to R/\mathfrak m$ is the original quotient map, and
$\rho:W_2(\kappa)\to\kappa$ is the original first-coordinate Witt
reduction. The square is an equality of entire original ring maps.
\end{theorem}
\begin{proof}
Choose $\mathfrak m$ and $p$ by the arithmetic closed-point theorem,
applied to this original $r$ and $N$. Thus $\kappa=R/\mathfrak m$ is a
finite field of characteristic $p$, $p>N$, and $\varphi(r)\ne0$.
A finite field is perfect: its Frobenius is injective and hence, on the
finite underlying set, bijective.

Use the original Witt reduction $\rho$, which is surjective, and put
$I=\ker\rho$. The proved original Witt arithmetic identifies every
$x\in I$ with a multiple $ap$ and proves $p^2=0$. Therefore for
$x=ap$ and $y=bp$ in $I$, the product is
$xy=abp^2=0$. Since $I^2$ is generated by these products, $I^2=0$.
This uses the actual original Witt arithmetic, without replacing the
base ring or assuming an infinitesimal-extension certificate.

Smoothness of $R$ over $\mathbb Z$ implies formal smoothness. Both
$\rho$ and $\varphi$ are actual $\mathbb Z$-algebra maps: every unital
ring map preserves the canonical integer maps. The infinitesimal lifting
property for the surjection $\rho$ with its square-zero kernel thus
constructs $\lambda:R\to W_2(\kappa)$ with
$\rho\lambda=\varphi$ as whole maps.

Finally, the quotient equivalence
$W_2(\kappa)/I\simeq\kappa$ carries the image of $\lambda(r)$ to
$\varphi(r)\ne0$, which is a unit in the original field. An element
whose image modulo a nilpotent ideal is a unit is itself a unit: lift an
inverse, so the product is $1+i$ with $i$ nilpotent, and invert $1+i$
by its finite geometric series. Apply this criterion to the square-zero
ideal $I$. Hence $\lambda(r)$ is a unit, and the constructed point
stays in the original principal open $D(r)$.
\end{proof}

This is a lift of the original arithmetic base point. It does not assert
descent of a family or line bundle, construct such a family, or prove a
Cartier, Deligne--Illusie, or vanishing statement.
